El sincrotró Alba és una instal·lació de recerca que utilitza llum generada per electrons accelerats per a analitzar les propietats i l'estructura de la matèria. Les principals qualitats d'aquesta radiació són un ampli espectre, una intensitat elevada i una brillantor extraordinària. Per a accelerar els electrons s'utilitzen camps elèctrics i magnètics. L'esquema mostra un model molt simplificat de funcionament: al començament del procés es generen electrons que s'acceleren en un accelerador lineal mitjançant un camp elèctric que suposarem uniforme al llarg de la zona d'acceleració, la qual té una longitud $d = 1,00\ \mathrm{m}$.

\figura[0.55]{fig1.png}

L'energia cinètica inicial dels electrons és zero, però quan surten de l'accelerador és d'1,00~keV.

\begin{apartat}{1}
Calculeu la intensitat del camp elèctric dins de l'accelerador i dibuixeu com són les línies de camp en aquesta regió.
\end{apartat}

\begin{apartat}{1}
Un cop els electrons han estat accelerats, se'ls condueix a l'anell de propulsió. Per a guiar els electrons al llarg de l'anell s'utilitzen camps magnètics. En l'esquema es mostra el primer camp magnètic que troben els electrons quan surten de l'accelerador lineal i entren a l'anell de propulsió. Si en aquesta regió no hi ha camp elèctric i el camp magnètic és de 0,15~T, calculeu la magnitud de la força que actuarà sobre l'electró. Quin tipus de trajectòria descriurà l'electró en aquesta regió? Justifiqueu la resposta.
\end{apartat}

\dades{%
  $m_\mathrm{e} = 9,11\times 10^{-31}\ \mathrm{kg}$.\\
  $|e| = 1,602\times 10^{-19}\ \mathrm{C}$.\\
  $1\ \mathrm{eV} = 1,602\times 10^{-19}\ \mathrm{J}$.}
